\section{Исходный код файла \texttt{MainWindow.cpp}}
\infoGen{Егоров Александр Владиславович}{swrneko@mai.education}
\footnotesize
\begin{listing}{1}
/****************************************************************
*                     КАФЕДРА № 304 2 КУРС                       *
*---------------------------------------------------------------*
* Project Type  : Qt Widgets Desktop Application                *
* Project Name  : FieldViz                                      *
* File Name     : MainWindow.cpp                                 *
* Language      : C/C++ (C++17)                                  *
* Programmer(s) : Егоров А.В. (swrneko)                          *
* Modifyed By   : Егоров А.В. (swrneko)                          *
* Edited by     : Neovim                                         *
* OS            : Arch Linux                                     *
* Created       : 01/07/26                                       *
* Last Revision : 01/07/26                                       *
* Comment(s)    : Реализация логики главного окна приложения.    *
****************************************************************/

#include "MainWindow.h"
#include <QCheckBox>
#include <QPushButton>
#include <QLabel>
#include <QVBoxLayout>
#include <QHBoxLayout>
#include <QGroupBox>

/****************************************
* Конструктор главного окна             *
****************************************/
MainWindow::MainWindow(QWidget *parent) : QMainWindow(parent) {
  setWindowTitle("КБН Портал: Визуализатор полей (Вариант 5)"); // Установка заголовка
  resize(850, 600); // Базовый размер окна

  setupUi(); // Инициализация интерфейса
}

/****************************************
* Инициализация элементов интерфейса    *
****************************************/
void MainWindow::setupUi() {
  QWidget *centralWidget = new QWidget(this); // Главный виджет
  setCentralWidget(centralWidget); // Установка

  QHBoxLayout *mainLayout = new QHBoxLayout(centralWidget); // Главный горизонтальный компоновщик

  m_visualizer = new Visualizer(this); // Создание 3D-сцены
  mainLayout->addWidget(m_visualizer, 3); // Добавление с весом 3

  QVBoxLayout *sidebar = new QVBoxLayout(); // Вертикальный компоновщик сайдбара
  mainLayout->addLayout(sidebar, 1); // Добавление с весом 1

  // Секция 1: Управление видимостью
  QGroupBox *grpControls = new QGroupBox("Отображение объектов", this); // Группа
  QVBoxLayout *ctrlLayout = new QVBoxLayout(grpControls); // Компоновщик группы

  m_chkFieldLines = new QCheckBox("Силовые линии поля a", this); // Чекбокс линий
  m_chkFieldLines->setChecked(true); // Включен
  connect(m_chkFieldLines, &QCheckBox::toggled, this, &MainWindow::onFieldLinesToggled); // Сигнал
  ctrlLayout->addWidget(m_chkFieldLines); // Добавление

  m_chkEquipotentials = new QCheckBox("Эквипотенциали поля b", this); // Чекбокс эквипотенциалей
  m_chkEquipotentials->setChecked(true); // Включен
  connect(m_chkEquipotentials, &QCheckBox::toggled, this, &MainWindow::onEquipotentialsToggled); // Сигнал
  ctrlLayout->addWidget(m_chkEquipotentials); // Добавление

  m_btnResetCamera = new QPushButton("Сбросить камеру", this); // Кнопка сброса
  connect(m_btnResetCamera, &QPushButton::clicked, this, &MainWindow::onResetCameraClicked); // Сигнал
  ctrlLayout->addWidget(m_btnResetCamera); // Добавление

  sidebar->addWidget(grpControls); // Добавление группы в сайдбар

  // Секция 2: Информационная панель (Solver)
  QGroupBox *grpInfo = new QGroupBox("Аналитический расчёт", this); // Группа
  QVBoxLayout *infoLayout = new QVBoxLayout(grpInfo); // Компоновщик

  m_lblInfoPanel = new QLabel(this); // Текстовая метка
  m_lblInfoPanel->setTextFormat(Qt::RichText); // Поддержка HTML
  m_lblInfoPanel->setWordWrap(true); // Перенос строк
  m_lblInfoPanel->setAlignment(Qt::AlignTop | Qt::AlignLeft); // Выравнивание

  m_lblInfoPanel->setText(
    "<b>Разработчик:</b><br>"
    "Егоров А.В. (М3О-225БВ-24)<br><br>"
    "<b>Индивидуальный вариант 5:</b><br>"
    "• Сфера S (1-й октант)<br><br>"
    "<b>Поле a (Соленоидальное):</b><br>"
    "<i>a = y²i + z²j + x²k</i><br>"
    "• div a = 0, rot a &ne; 0<br>"
    "• Поток P = 3&pi;/16 &approx; 0.5890<br>"
    "• Циркуляция C = -2<br><br>"
    "<b>Поле b (Потенциальное):</b><br>"
    "<i>b = (2xy³z+y)i + ...</i><br>"
    "• rot b = 0<br>"
    "• Потенциал U = xy + x²y³z + z⁴<br>"
    "• Работа W (O &rarr; M) = 3"
  ); // HTML Текст

  infoLayout->addWidget(m_lblInfoPanel); // Добавление в группу
  sidebar->addWidget(grpInfo); // Добавление группы в сайдбар

  sidebar->addStretch(); // Растяжка для прижатия кверху
}

/****************************************
* Обработчики сигналов управления       *
****************************************/
void MainWindow::onFieldLinesToggled(bool checked) {
  m_visualizer->setShowFieldLines(checked); // Переключение линий
}

void MainWindow::onEquipotentialsToggled(bool checked) {
  m_visualizer->setShowEquipotentials(checked); // Переключение эквипотенциалей
}

void MainWindow::onResetCameraClicked() {
  m_visualizer->resetCamera(); // Сброс камеры
}

\end{listing}
\normalsize
\pagebreak
